\chapter{全存储顺序以及 x86 内存模型}

\begin{epigraph}
\textbf{核心命题}　x86 采用的全存储顺序（TSO）在 SC 与松弛模型之间取折中——唯一允许的重排是 store$\rightarrow$load，这让 store buffer 可以让后续 load 绕过未完成的 store。Sewell 等人 2010 年的 x86-TSO 形式化终结了 Intel 长达数年的"processor ordering"歧义。
\end{epigraph}



\section{TSO/x86 的动机}

SC 的根本性能瓶颈在 store：store 必须等其他核的 invalidate ack 全部返回才能提交。但现实代码中，store$\rightarrow$load 是最常见的相邻模式——例如先写 flag 再读共享数据，或先写参数再调用函数。如果能让 load 绕过自己未完成的 store，就能大幅隐藏 store 延迟。

这正是 \textbf{TSO（Total Store Order）}的动机：在 SC 基础上\textbf{唯一放松 store$\rightarrow$load 重排}，让 store buffer 可以让后续 load 先行（从 buffer forwarding 或从内存读）。

TSO 最早由 SPARC 在 1992 年的 SPARC V8 架构手册中正式定义（"Total Store Ordering"）。x86 实质采用 TSO，但 Intel 在早期文档中用模糊的"processor ordering"描述，直到 2007–2008 年才正式澄清为 TSO。

\section{TSO/x86 的基本思想}

TSO 的核心机制是\textbf{每核一个 FIFO store buffer}：

\begin{itemize}
\item store 先进入本核 store buffer 的尾部，按 FIFO 顺序排入内存；
\item load 时，若本核 store buffer 中有同地址的最近 store，必须先读它（store forwarding）；
\item 否则，load 从内存读。
\end{itemize}


这个机制的关键后果：\textbf{本核的 load 可以"跳过"自己未排入内存的 store}——因为 load 可能从内存读到旧值（如果 store 还在 buffer 里）。这就是 store$\rightarrow$load 重排的物理来源。

\section{TSO/x86 的形式化描述}

Sewell、Sarkar、Owens 等人在《x86-TSO: A Rigorous and Usable Programmer's Model for x86 Multiprocessors》（\emph{Communications of the ACM} 53(7): 89–97, 2010）中给出了 TSO 的严格形式化。核心要素：

\begin{itemize}
\item 每个线程与共享内存之间插入一个\textbf{无界 FIFO store buffer}；
\item 线程读时，若自己 buffer 中有同地址的最近 store，必须先读它（store forwarding）；
\item 否则从共享内存读；
\item store 按全局 FIFO 总序（total store order）排入内存。
\end{itemize}


这个形式化有两种等价表达：\textbf{操作化模型}（抽象机，store buffer + 全局锁）与\textbf{公理化模型}（happens-before + acyclicity）。Owens/Sarkar/Sewell 在 TPHOLs 2009 给出了两种模型的等价证明（在 HOL 定理证明器中机械化）。

\section{实现 TSO/x86}

TSO 的实现相对简单：

\begin{itemize}
\item \textbf{硬件}：每核一个 FIFO store buffer；store 进入 buffer 尾部，按 FIFO drain 到内存；
\item \textbf{load 路径}：先查 store buffer（同地址 forwarding），未命中再查缓存/内存；
\item \textbf{store 路径}：进 buffer，异步 drain。
\end{itemize}


这个实现的关键优势是 store 不再需要等 invalidate ack 才能让后续 load 进行——store buffer 吸收了 store 延迟。

\subsection{实现原子指令}

x86 的原子指令通过 \textbf{LOCK 前缀}实现（如 \texttt{lock cmpxchg}、\texttt{lock xchg}、\texttt{lock add}）：

\begin{itemize}
\item LOCK 前缀在 drain store buffer 的同时提供原子性；
\item 行为等价于全 fence（但额外保证 read-modify-write 的原子性）；
\item 常见技巧：用 \texttt{lock addq \$0,(\%rsp)}——一条带 LOCK 前缀的无关 store 充当便宜的 barrier。
\end{itemize}


\subsection{实现 FENCE}

\texttt{MFENCE} 指令是全屏障：串行化此前所有 load/store（drain store buffer 与 load buffer），阻止 store$\rightarrow$load 重排，恢复 SC 语义。在需要 SC 语义的关键段（如 Dekker 算法、双重检查锁）必须插入 MFENCE。

\section{关于 TSO 的进一步阅读资料}

\begin{itemize}
\item SPARC V8 架构手册（1992）：TSO 的最早正式定义；
\item Sewell et al. CACM 2010：x86-TSO 的严格形式化；
\item Owens/Sarkar/Sewell TPHOLs 2009：x86-TSO 的双模型等价证明；
\item Intel SDM Vol.3 §8.2：x86 内存排序的官方规范。
\end{itemize}


\section{比较 SC 和 TSO}

SC 与 TSO 的关键差异：

\begin{center}
\begin{tabular}{lll}
\toprule
特性 & SC & TSO \\
\midrule
store$\rightarrow$load 重排 & 禁止 & \textbf{允许} \\
store$\rightarrow$store 重排 & 禁止 & 禁止（store 有全局总序） \\
load$\rightarrow$load 重排 & 禁止 & 禁止 \\
load$\rightarrow$store 重排 & 禁止 & 禁止 \\
store buffer & 受限 & 充分利用 \\
典型实现 & MIPS R10000 & x86、SPARC V8 \\
\bottomrule
\end{tabular}
\end{center}


唯一差异就是 store$\rightarrow$load 重排。这看似细微，但对性能影响巨大——它让 store buffer 可以满负荷工作。代价是程序员在某些场景（如 Dekker 算法、发布模式）必须显式插入 MFENCE。

\section{小结}

TSO 是 x86 与 SPARC V8 采用的内存模型，唯一放松 SC 的是 store$\rightarrow$load 重排。Sewell 等人 2010 年的 x86-TSO 形式化终结了 Intel "processor ordering" 的歧义。TSO 的实现是每核 FIFO store buffer + 全局 store 总序，store forwarding 保证 load 能读到自己最近的 store。原子指令通过 LOCK 前缀实现，MFENCE 恢复 SC 语义。SC 与 TSO 的唯一差异是 store$\rightarrow$load 重排，但这对性能影响巨大。

\section*{参考文献}
\begin{itemize}
\item Sewell P., Sarkar S., Owens S., et al. x86-TSO: A rigorous and usable programmer's model for x86 multiprocessors. \emph{Communications of the ACM}, 2010, 53(7): 89–97. https://doi.org/10.1145/1785414.1785443
\item Owens S., Sarkar S., Sewell P. A better x86 memory model: x86-TSO. TPHOLs 2009. https://www.cl.cam.ac.uk/\textasciitilde{}pes20/weakmemory/x86tso-paper.tphols.pdf
\item SPARC International. \emph{The SPARC Architecture Manual (Version 8)}. 1992. https://courses.grainger.illinois.edu/cs423/sp2011/lectures/sim\_public/sparcv8.pdf
\item Intel. \emph{Intel 64 Architecture Memory Ordering White Paper}. Document 318147, 2007. https://www.cs.cmu.edu/\textasciitilde{}410-f10/doc/Intel\_Reordering\_318147.pdf
\item Intel. \emph{Intel 64 and IA-32 Architectures Software Developer's Manual, Vol. 3, §8.2 Memory Ordering}. https://www.intel.com/content/www/us/en/developer/articles/technical/intel-sdm.html
\item Sorin D. J., Hill M. D., Wood D. A. \emph{A Primer on Memory Consistency and Cache Coherence} (2nd ed.). 2020, Chapter 4. https://pages.cs.wisc.edu/\textasciitilde{}markhill/papers/primer2020\_2nd\_edition.pdf
\item Higham L., Kaverheid L. Memory consistency and process coordination for SPARC. Springer. https://link.springer.com/chapter/10.1007/3-540-44467-x\_32
\end{itemize}
